\part{字符串}

\chapter{字符串}

\section{哈希}
\entrymeta{$O(1)$ 单次查询}
使用 unsigned long long 的模 2\^{}64 自然溢出；反转等价类保存排序后的正反哈希二元组， 避免把两个哈希相乘导致的额外信息丢失。哈希仍有碰撞可能，不提供确定性判等保证。

\subsection{P3498 {[}POI 2010{]} KOR-Beads}\label{p3498-poi-2010-kor-beads}






\begin{codeline}
#define N 200005
#define BASE 200019
class DeHash {
public:
    uint64 asc, rev;
};
inline bool operator<(const DeHash& A, const DeHash& B) {
    if(A.asc != B.asc) {
        return A.asc < B.asc;
    } else {
        return A.rev < B.rev;
    }
}
inline bool operator==(const DeHash& A, const DeHash& B) {
    return A.asc == B.asc && A.rev == B.rev || A.asc == B.rev && A.rev == B.asc;
}
DeHash de_hash[N];
set<pair<uint64, uint64>> s;
uint64 base[N];
uint64 n;
uint64 bead[N];
/*
 * 返回闭区间[l,r]子串的正反向哈希，反转后的字符串视为同类
 */
inline DeHash getSegHash(uint64 l, uint64 r) {
    return {
        de_hash[r].asc - de_hash[l - 1].asc * base[r - l + 1],
        de_hash[l].rev - de_hash[r + 1].rev * base[r - l + 1]};
}

int main() {
    ios::sync_with_stdio(false);
    cin.tie(nullptr);
    cin >> n;
    for(uint64 i = 1; i <= n; i++) {
        cin >> bead[i];
    }
    base[0] = 1;
    for(uint64 i = 1; i <= n; i++) {
        base[i] = base[i - 1] * BASE;
        de_hash[i].asc = de_hash[i - 1].asc * BASE + bead[i];
        de_hash[n - i + 1].rev = de_hash[n - i + 2].rev * BASE + bead[n - i + 1];
    }
    vector<uint64> ans_k;
    uint64 max_seg = 0, size;
    DeHash tmp;
\end{codeline}
\vspace{-0.9\baselineskip}
\begin{codeline}[firstnumber=47]
    for(uint64 i = 1; i <= n; i++) {
        s.clear();

        for(uint64 j = 1; j + i - 1 <= n; j += i) {
            tmp = getSegHash(j, j + i - 1);
            if(!s.count(minmax(tmp.asc, tmp.rev))) {
                s.insert(minmax(tmp.asc, tmp.rev));
            } else {
            }
        }
        size = s.size();
        if(size > max_seg) {
            max_seg = size;
            ans_k.clear();
            ans_k.push_back(i);
        } else if(size == max_seg) {
            ans_k.push_back(i);
        }
    }
    cout << (max_seg) << ' ' << (ans_k.size()) << '\n';
    for(uint64 e : ans_k) {
        cout << (e) << ' ';
    }
}
\end{codeline}
\begin{warnbox}
双哈希降低碰撞；模数与底数取大质数。
\end{warnbox}

\section{KMP}
\entrymeta{$O(n+m)$}
\subsection{P3375 【模板】KMP}\label{p3375-ux6a21ux677fkmp}


\begin{codeline}
class KMP {
public:
    int n;
    string pat;
    vector<int> nxt;
    /*
     * 由模式串pat预处理失配数组
     */
    explicit KMP(const string& pat) : n(int(pat.size())), pat(pat) {
        nxt.resize(n + 1);
        nxt[0] = -1;
        for(int i = 0, j = -1; i < n;) {
            if(j == -1 || pat[i] == pat[j]) {
                ++i;
                ++j;
                nxt[i] = j;
            } else {
                j = nxt[j];
            }
        }
    }

    /*
     * 返回模式串在s中出现的所有起点，下标从0开始
     */
    vector<int> search(const string& s) const {
        vector<int> ans;
        int m = s.size();
        if(n == 0) {
            for(int i = 0; i <= m; i++) {
                ans.push_back(i);
            }
            return ans;
        }
        for(int i = 0, j = 0; i < m;) {
            if(j == -1 || s[i] == pat[j]) {
                ++i;
                ++j;
                if(j == n) {
                    ans.push_back(i - j);
                    j = nxt[j];
                }
            } else {
                j = nxt[j];
            }
        }
\end{codeline}
\vspace{-0.9\baselineskip}
\begin{codeline}[firstnumber=47]
        return ans;
    }
};
\end{codeline}


\subsection{P9770 {[}HUSTFC 2023{]} 逆 KMP}\label{p9770-hustfc-2023-ux9006-kmp}


\begin{codeline}
const int maxn = 200005;
int ans[maxn], n, fa[maxn][18];
int find(int x, int k) {
    return fa[x][k] == x ? x : fa[x][k] = find(fa[x][k], k);
}
void merge(int x, int y, int k) {
    x = find(x, k);
    y = find(y, k);
    if(x < y) {
        swap(x, y);
    }
    if(x != y) {
        fa[x][k] = y;
    }
}
int main() {
    ios::sync_with_stdio(false);
    cin.tie(nullptr);
    cin >> n;
    const int maxk = __lg(n);
    for(int i = 1; i <= n; ++i) {
        for(int k = 0; k <= maxk; ++k) {
            fa[i][k] = i;
        }
    }
    for(int i = 1, l1, r1, l2, r2; i <= n; ++i) {
        int x;
        cin >> x;
        if(x) {
            l1 = 1;
            r1 = x;
            l2 = i - x + 1;
            r2 = i;
            for(int k = maxk; ~k; --k) {
                if(l1 + (1 << k) - 1 <= r1) {
                    merge(l2, l1, k);
                    l1 += 1 << k;
                    l2 += 1 << k;
                }
            }
        }
    }
    for(int k = maxk; k; --k) {
        for(int i = 1; i + (1 << k) - 1 <= n; ++i) {
            int pos = find(i, k);
            merge(i, pos, k - 1);
\end{codeline}
\vspace{-0.9\baselineskip}
\begin{codeline}[firstnumber=47]
            merge(i + (1 << (k - 1)), pos + (1 << (k - 1)), k - 1);
        }
    }
    int cnt = 0;
    for(int i = 1; i <= n; ++i) {
        if(find(i, 0) == i) {
            ans[i] = ++cnt;
        } else {
            ans[i] = ans[find(i, 0)];
        }
    }
    for(int i = 1; i <= n; ++i) {
        cout << (ans[i]) << ' ';
    }
    return 0;
}
\end{codeline}
\begin{warnbox}
\texttt{nxt[0] = -1}；匹配成功后 \texttt{j = nxt[j]} 继续找重叠匹配。
\end{warnbox}

\section{AC 自动机}
\entrymeta{$O(\sum|s|)$ 建树}
\subsection{P5357 【模板】AC 自动机}\label{p5357-ux6a21ux677fac-ux81eaux52a8ux673a}


\begin{codeline}
class ACAutomaton {
public:
    class Node {
    public:
        int next[26]{}, fail = 0;
    };
    vector<Node> node{Node()};
    vector<int> order;
    bool built = false;
    /*
     * 清空自动机
     */
    void clear() {
        node.assign(1, Node());
        order.clear();
        built = false;
    }
    /*
     * 插入模式串s，返回终止节点编号
     */
    int insert(const string& s) {
        int u = 0;
        for(char c : s) {
            int x = c - 'a';

            if(!node[u].next[x]) {
                int v = int(node.size());
                node.emplace_back();
                node[u].next[x] = v;
            }
            u = node[u].next[x];
        }
        return u;
    }
    /*
     * 插入全部模式串后建立fail指针和转移
     */
    void build() {
        if(built) {
            return;
        }
        for(int v : node[0].next) {
            if(v) {
                order.push_back(v);
            }
        }
\end{codeline}
\vspace{-0.9\baselineskip}
\begin{codeline}[firstnumber=47]
        for(int i = 0; i < int(order.size()); i++) {
            int u = order[i];
            for(int c = 0; c < 26; c++) {
                int v = node[u].next[c];
                if(v) {
                    node[v].fail = node[node[u].fail].next[c];
                    order.push_back(v);
                } else {
                    node[u].next[c] = node[node[u].fail].next[c];
                }
            }
        }
        built = true;
    }
    /*
     * 统计文本s中各节点对应模式串的出现次数，重复查询不累积
     */
    vector<long long> query(const string& s) const {
        vector<long long> ans(node.size());
        int u = 0;
        for(char c : s) {
            u = node[u].next[c - 'a'];
            ans[u]++;
        }
        for(auto it = order.rbegin(); it != order.rend(); ++it) {
            ans[node[*it].fail] += ans[*it];
        }
        return ans;
    }
};
\end{codeline}
\begin{warnbox}
先全部插入再 \texttt{build()}；fail 树上逆 BFS 累加。
\end{warnbox}

\section{后缀自动机}
\entrymeta{$O(n)$}
\subsection{题意翻译}\label{ux9898ux610fux7ffbux8bd1}

输入两个字符串，输出它们的最长公共子串长度，若不存在公共子串则输出 0。


\begin{codeline}
class State {
public:
    int len{0}, link{0};
    map<char, int> next;
};

class SuffixAutomaton {
    vector<State> state;
    int size{0}, last{0};

public:
    explicit SuffixAutomaton(int n = 0) {
        init(n);
    }

    /*
     * 清空自动机，n为预计字符串长度
     */
    void init(int n = 0) {
        state.assign(1, State());
        state.reserve(max(1, 2 * n));
        state[0].link = -1;
        size = 1;
        last = 0;
    }

    /*
     * 在当前字符串末尾加入字符c
     */
    void extend(char c) {
        int cur = size++;
        state.emplace_back();
        state[cur].len = state[last].len + 1;
        int p = last;
        while(p != -1 && !state[p].next.count(c)) {
            state[p].next[c] = cur;
            p = state[p].link;
        }
        if(p == -1) {
            state[cur].link = 0;
        } else {
            int q = state[p].next[c];
            if(state[p].len + 1 == state[q].len) {
                state[cur].link = q;
            } else {
                int clone = size++;
\end{codeline}
\vspace{-0.9\baselineskip}
\begin{codeline}[firstnumber=47]
                state.emplace_back();
                state[clone].len = state[p].len + 1;
                state[clone].next = state[q].next;
                state[clone].link = state[q].link;
                while(p != -1 && state[p].next[c] == q) {
                    state[p].next[c] = clone;
                    p = state[p].link;
                }
                state[q].link = state[cur].link = clone;
            }
        }
        last = cur;
    }

    /*
     * 返回两个字符串的一个最长公共子串，调用时重建自动机
     */
    string lcs(const string& string1, const string& string2) {
        init(int(string1.size()));
        for(char c : string1) {
            extend(c);
        }
        int v = 0, l = 0, best = 0, bestpos = -1;
        for(int i = 0; i < string2.size(); i++) {
            while(v && !state[v].next.count(string2[i])) {
                v = state[v].link;
                l = state[v].len;
            }
            if(state[v].next.count(string2[i])) {
                v = state[v].next[string2[i]];
                l++;
            }
            if(l > best) {
                best = l;
                bestpos = i;
            }
        }
        return string2.substr(bestpos - best + 1, best);
    }

};
\end{codeline}


\href{../examples/字符串/后缀自动机/01.cpp}{题目输入输出示例}（与正文组装后编译，正文不重复驱动代码）。

\begin{warnbox}
clone 节点完整复制转移与 link，再重定向。
\end{warnbox}

